% =====================================================================
% SECCIÓN 2: MARCO METODOLÓGICO, OBJETIVOS E HIPÓTESIS
% =====================================================================

% ---------------------------------------------------------------------
% Diapositiva 4: Objetivos
% ---------------------------------------------------------------------
\begin{frame}{Objetivos de la investigación}
  \framesubtitle{Estructura articulada según la taxonomía de Bloom}

  \ugrblock{Objetivo general (Bloom 6 -- Crear)}{
    {\footnotesize Desarrollar un modelo de aprendizaje automático supervisado que, a partir de texto escrito en español, clasifique publicaciones de usuarios con indicios de sintomatología depresiva mediante un modelo de lenguaje contextual.}
  }

  \begin{columns}[T]
    \column{0.49\textwidth}
    \ugrblock{OE1 y OE2: Datos y líneas base}{
      {\scriptsize
      \begin{itemize}
        \item \textbf{OE1 -- Consolidar:} Integrar y estructurar un dataset masivo en español a partir de redes sociales.
        \item \textbf{OE2 -- Entrenar:} Entrenar clasificadores supervisados basados en representaciones léxicas TF-IDF.
      \end{itemize}}
    }
    \column{0.49\textwidth}
    \ugrblock{OE3 -- OE5: Contexto y evaluación}{
      {\scriptsize
      \begin{itemize}
        \item \textbf{OE3 -- Fine-tuning:} Ajuste fino supervisado de BETO.
        \item \textbf{OE4 -- Comparar:} Contraste empírico multidimensional.
        \item \textbf{OE5 -- Seleccionar:} Selección según significación ($p < 0{,}0005$), unidad de análisis y costo.
      \end{itemize}}
    }
  \end{columns}

  \note{
    Se definieron cinco objetivos específicos alineados con la taxonomía de Bloom: consolidar el corpus masivo integrado (OE1), entrenar líneas base léxicas (OE2), realizar el fine-tuning de BETO (OE3), comparar exhaustivamente los modelos (OE4) y seleccionar la arquitectura óptima según métricas y significación estadística (OE5).
  }
\end{frame}

% ---------------------------------------------------------------------
% Diapositiva 5: Hipótesis de Investigación y Protocolo Formal
% ---------------------------------------------------------------------
\begin{frame}{Hipótesis formal y protocolo de contraste}
  \framesubtitle{Criterio estadístico riguroso de decisión formal}

  \begin{columns}[T]
    \column{0.49\textwidth}
    \ugrblockok{Hipótesis de investigación ($H_1$)}{
      {\scriptsize Es posible desarrollar un modelo de aprendizaje automático supervisado que, a partir de texto escrito en español, clasifique publicaciones de usuarios con indicios de sintomatología depresiva mediante un modelo de lenguaje contextual.}
    }
    \column{0.49\textwidth}
    \ugrblock{Hipótesis nula ($H_0$)}{
      {\scriptsize No es posible desarrollar un modelo de aprendizaje automático supervisado que, a partir de texto escrito en español, clasifique publicaciones de usuarios con indicios de sintomatología depresiva mediante un modelo de lenguaje contextual.}
    }
  \end{columns}

  \vskip-0.05cm
  \ugrblock{Criterio de decisión formal (Bootstrap por conglomerados de usuarios)}{
    {\scriptsize Se define el contraste formal con $B = 2\,000$ réplicas remuestreando los 64 autores independientes del test nativo ($\alpha = 0{,}05$):
    \begin{equation*}
      \text{Rechazar } H_0 \text{ a favor de } H_1 \iff \Delta F_{1\text{-macro}}^{\text{low}} > 0 \;\land\; \Delta\text{AUC-ROC}^{\text{low}} > 0
    \end{equation*}
    donde $\Delta = \text{Métrica}(\text{BETO}_{\text{mono}}) - \text{Métrica}(\text{Mejor Línea Base})$. Los extremos inferiores del IC 95\,\% pareado deben excluir estrictamente el valor nulo ($0$).}
  }

  \note{
    La hipótesis se presenta verbatim conforme al informe final. La condición formal de rechazo de H0 exige que los extremos inferiores de los intervalos de confianza del 95\% para la diferencia pareada de F1-macro y AUC-ROC sean estrictamente positivos en el remuestreo de bootstrap por conglomerados de usuarios.
  }
\end{frame}

% ---------------------------------------------------------------------
% Diapositiva 6: Metodología CRISP-DM
% ---------------------------------------------------------------------
\begin{frame}{Metodología de trabajo: Enfoque CRISP-DM}
  \framesubtitle{Proceso estándar de minería de datos articulado en seis fases iterativas}

  \vskip-0.2cm
  \begin{center}
    \resizebox{0.88\textwidth}{!}{% Proceso de investigación según CRISP-DM (seis fases, recorrido en serpentina)
% Fila superior: 1 -> 2 -> 3. Fila inferior: 4 -> 5 -> 6 (de derecha a izquierda).
% Requiere los estilos proceso, resultado y flecha del documento anfitrión.
\begin{tikzpicture}[font=\footnotesize]
    \hyphenpenalty=10000\relax

    \tikzset{
        fase/.style={proceso=4.9cm, minimum height=3.0cm, align=left, anchor=center},
        iter/.style={-{Stealth[length=2.4mm, width=1.8mm]}, line width=0.8pt, color=ugrNaranja, dashed},
        ida/.style={{Stealth[length=2.4mm, width=1.8mm]}-{Stealth[length=2.4mm, width=1.8mm]}, line width=0.8pt, color=ugrGris},
    }

    \node[fase] (f1) at (-6.0, 0) {%
        {\small\bfseries\color{ugrAzul} 1. Comprensión del problema}\\[2pt]
        Pregunta, $H_1$/$H_0$ y objetivos\\
        Criterio de éxito: IC 95\,\% de $\Delta F_1$ y $\Delta$AUC $>0$\\
        Restricciones: tarea binaria, GPU de 8\,GB, corpus abiertos};

    \node[fase] (f2) at (0, 0) {%
        {\small\bfseries\color{ugrAzul} 2. Comprensión de los datos}\\[2pt]
        Cuatro fuentes: Twitter nativo y tres de Reddit traducido\\
        Análisis exploratorio, calidad y desbalance\\
        Distribución de clases y longitudes};

    \node[fase] (f3) at (6.0, 0) {%
        {\small\bfseries\color{ugrAzul} 3. Preparación de los datos}\\[2pt]
        Anonimización y traducción (MarianMT)\\
        Corpus unificado (1\,311\,023 textos)\\
        Partición 70/10/20 por usuario, rasgos y balanceo};

    \node[fase, draw=ugrAzul, fill=ugrAzul!8] (f4) at (6.0, -4.6) {%
        {\small\bfseries\color{ugrAzul} 4. Modelado}\\[2pt]
        TF-IDF con Regresión Logística y SVM\\
        Modelo combinado (19 rasgos)\\
        BETO; esquemas mono y multi};

    \node[fase, draw=ugrVerde, fill=ugrVerde!9] (f5) at (0, -4.6) {%
        {\small\bfseries\color{ugrAzul} 5. Evaluación}\\[2pt]
        $F_1$-macro, AUC, $\kappa$ y \emph{bootstrap} de 2\,000 réplicas\\
        Nivel de usuario y contraste formal\\
        Contraste de $H_0$};

    \node[fase] (f6) at (-6.0, -4.6) {%
        {\small\bfseries\color{ugrAzul} 6. Despliegue y cierre}\\[2pt]
        Código, manifiestos y predicciones reproducibles\\
        Uso no clínico y consideraciones éticas\\
        Informe y presentación};

    \draw[ida]   (f1) -- (f2);
    \draw[ida]   (f2) -- (f3);
    \draw[ida]   (f3) -- (f4);
    \draw[flecha] (f4) -- (f5);
    \draw[flecha] (f5) -- (f6);
    \draw[iter]  (f5.north west) -- node[midway, right=3pt, font=\scriptsize, text=ugrNaranja, align=left, inner sep=1pt] {nuevo\\ciclo} (f1.south east);

\end{tikzpicture}
}
  \end{center}
  \vskip-0.35cm
  {\scriptsize\centering\color{ugrGris}El flujo iterativo conecta la evaluación con la comprensión de datos y el modelado; el repositorio garantiza reproducibilidad con \texttt{make}.\par}

  \note{
    Adoptamos el estándar de Ciencia de Datos CRISP-DM en seis fases interconectadas. El ciclo conecta permanentemente la evaluación analítica con el preprocesamiento y el modelado. Todo el pipeline es reproducible mediante comandos en el Makefile del repositorio.
  }
\end{frame}
